\begin{definitionbox}{Definition (Left Inverse)}
\begin{definition}[Left Inverse]
\label{def:left-inverse}
An element $x'$ is a left inverse of $x$ with respect to identity $e$ if
\[
x'\mathbin{O}x=e.
\]
\end{definition}
\end{definitionbox}

\begin{definitionbox}{Definition (Right Inverse)}
\begin{definition}[Right Inverse]
\label{def:right-inverse}
An element $x''$ is a right inverse of $x$ with respect to identity $e$ if
\[
x\mathbin{O}x''=e.
\]
\end{definition}
\end{definitionbox}

\begin{definitionbox}{Definition (Two-Sided Inverse)}
\begin{definition}[Two-Sided Inverse]
\label{def:two-sided-inverse}
An element $y$ is a two-sided inverse of $x$ if
\[
y\mathbin{O}x=e
\quad\text{and}\quad
x\mathbin{O}y=e.
\]
\end{definition}
\end{definitionbox}

\begin{lemmabox}{Lemma (Left and Right Inverses Coincide)}
\begin{lemma}[Left and Right Inverses Coincide]
\label{lem:left-right-inverses-coincide}
In an associative operation with identity, a left inverse and a right inverse
of the same element are equal.
\end{lemma}
\end{lemmabox}

\begin{remark*}[Standard quantified statement]
\[
(x'\mathbin{O}x=e)\land(x\mathbin{O}x''=e)
\Longrightarrow x'=x''.
\]
\end{remark*}

\begin{dependencies}
\begin{itemize}
  \item \hyperref[def:associativity]{Associativity}
  \item \hyperref[def:two-sided-identity]{Two-Sided Identity}
  \item \hyperref[def:left-inverse]{Left Inverse}
  \item \hyperref[def:right-inverse]{Right Inverse}
\end{itemize}
\end{dependencies}

\begin{theorembox}{Theorem (Uniqueness of Inverse)}
\begin{theorem}[Uniqueness of Inverse]
\label{thm:uniqueness-of-inverse}
In an associative operation with identity, inverses are unique.
\end{theorem}
\end{theorembox}

\begin{dependencies}
\begin{itemize}
  \item \hyperref[lem:left-right-inverses-coincide]{Left and Right Inverses Coincide}
\end{itemize}
\end{dependencies}

\begin{definitionbox}{Definition (Right Solvability)}
\begin{definition}[Right Solvability]
\label{def:right-solvability}
Right solvability means
\[
\forall x,y\,\exists z,\ x=y\mathbin{O}z,
\]
\end{definition}
\end{definitionbox}

\begin{definitionbox}{Definition (Left Solvability)}
\begin{definition}[Left Solvability]
\label{def:left-solvability}
Left solvability means
\[
\forall x,y\,\exists z,\ x=z\mathbin{O}y.
\]
\end{definition}
\end{definitionbox}

\begin{definitionbox}{Definition (Two-Sided Solvability)}
\begin{definition}[Two-Sided Solvability]
\label{def:two-sided-solvability}
Two-sided solvability means both left solvability and right solvability hold.
\end{definition}
\end{definitionbox}

\begin{propositionbox}{Proposition (Solvability and Inverses)}
\begin{proposition}[Solvability and Inverses]
\label{prop:solvability-iff-inverses}
In a monoid, solvability is equivalent to the existence of inverses.
\end{proposition}
\end{propositionbox}

\begin{dependencies}
\begin{itemize}
  \item \hyperref[def:monoid]{Monoid}
  \item \hyperref[def:two-sided-inverse]{Two-Sided Inverse}
  \item \hyperref[def:right-solvability]{Right Solvability}
  \item \hyperref[def:left-solvability]{Left Solvability}
  \item \hyperref[def:two-sided-solvability]{Two-Sided Solvability}
\end{itemize}
\end{dependencies}
